% Lecture notes: Week 7, Lecture 3 -- Beam Deflection
% To hide this lecture from the website, add a line containing only:  % draft
% To set the date shown on the website, add a line like:  % posted: 2026-09-02
\documentclass[11pt]{article}
\usepackage{mech2030notes}
\begin{document}

% ##################################################################
%  WEEK 7, LECTURE 3
% ##################################################################
\lecture{7}{3}{Beam Deflection}

\section*{Notation}

\begin{center}
\begin{tikzpicture}
  % deflected shape (drawn upward, the positive direction):  v = 1.1 (x/L)^2 (3 - x/L)/2
  \pgfmathsetmacro{\Lb}{6}
  \pgfmathsetmacro{\xs}{4.4}
  \pgfmathsetmacro{\vs}{1.6*(\xs/\Lb)^2*(3-\xs/\Lb)/2}
  \pgfmathsetmacro{\ss}{1.6*(6*\xs/\Lb - 3*(\xs/\Lb)^2)/(2*\Lb)}
  \pgfmathsetmacro{\ang}{atan(\ss)}
  \pgfmathsetmacro{\xd}{2.6}
  \pgfmathsetmacro{\vd}{1.6*(\xd/\Lb)^2*(3-\xd/\Lb)/2}
  \lwall{0}{-0.9}{2.0}
  \draw[thin, dashed] (0,0) -- (6.6,0);
  \draw[plotline] plot[domain=0:\Lb, samples=50] (\x,{1.6*(\x/\Lb)^2*(3-\x/\Lb)/2});
  % deflection v(x)
  \draw[force] (\xd,0) -- (\xd,\vd);
  \node[right] at (\xd,{\vd/2+0.05}) {$v(x)$};
  % slope theta(x)
  \draw[thin] ({\xs-1.0},{\vs-1.0*\ss}) -- ({\xs+1.6},{\vs+1.6*\ss});
  \draw[thin, dashed] (\xs,\vs) -- ({\xs+1.9},\vs);
  \draw ({\xs+1.4},\vs) arc (0:\ang:1.4);
  \node[right] at ({\xs+1.45},{\vs+0.2}) {$\theta(x)$};
  \node[pin] at (\xs,\vs) {};
  % x axis
  \draw[dimto] (0,-1.3) -- (1.0,-1.3) node[right] {$x$};
\end{tikzpicture}
\end{center}

\begin{itemize}
  \item $v(x)$ is the \emph{deflection} of the beam (positive \emph{upwards}).
  \item $\theta(x)$ is the \emph{slope angle} of the beam, where, since slopes are small,
        \[ \boxed{\theta(x) \approx \tan\theta = \frac{dv}{dx}} \]
  \item The \emph{curvature} of the beam is
        \[ \boxed{\kappa = \frac{d^2 v}{dx^2}} \]
\end{itemize}

\section*{Moment--Curvature Relationship}

\[
  \boxed{M = \kappa\, EI} \qquad \text{(derivation in Hibbeler, Ch.~12)}
\]

\noindent
$\hookrightarrow$ This means
\begin{align*}
  & \boxed{M = \frac{d^2 v}{dx^2}\,EI} \\[6pt]
  \text{or} \quad & \frac{dM}{dx} = \frac{d^3 v}{dx^3}\,EI = V(x)
     && \left(\text{since } \frac{dM}{dx} = V(x)\right) \\[6pt]
  \text{or} \quad & \frac{dV}{dx} = \frac{d^4 v}{dx^4}\,EI = p(x)
     && \left(\text{since } \frac{dV}{dx} = p(x)\right)
\end{align*}

\noindent
$\rightarrow$ This is the \emph{Euler--Bernoulli beam equation}:
\begin{center}
\fbox{\fbox{\parbox{0.42\textwidth}{\centering
  \[ EI\,\frac{d^4 v}{dx^4} = p(x) \]}}}\\[0.4em]
(Assumes constant $EI$ and small deflections.)
\end{center}

\subsection*{Notes}

\noindent
Each derivative has a physical meaning:
\[
  \left\{
  \begin{aligned}
    p(x) &= \frac{d^4 v}{dx^4}\,EI && \text{(distributed load)} \\[2pt]
    V(x) &= \frac{d^3 v}{dx^3}\,EI && \text{(shear force)} \\[2pt]
    M(x) &= \frac{d^2 v}{dx^2}\,EI && \text{(bending moment)} \\[2pt]
    \theta(x) &= \frac{dv}{dx} && \text{(slope)} \\[2pt]
    v(x) & && \text{(deflection)}
  \end{aligned}
  \right.
\]

\begin{itemize}[label=\then]
  \item The solution requires \emph{integrating four times}.
  \item We need \emph{four boundary conditions}.
\end{itemize}

% ==================================================================
\section*{Example: Squirrel Problem}
% ==================================================================

\noindent
\begin{minipage}[c]{0.64\textwidth}
\centering
\begin{tikzpicture}
  \lwall{0}{-1.5}{1.0}
  \draw[beam] (0,-0.12) rectangle (5,0.12);
  \draw[thin, dashed] (0,0) -- (5,0);
  \draw[plotline] plot[domain=0:5, samples=40] (\x,{-1.1*(\x/5)^2*(3-\x/5)/2});
  \node[pin] at (5,-1.1) {};
  \draw[force] (5,1.3) node[above] {$P = \SI{5}{N}$} -- (5,0.14);
  \draw[force] (5.6,0) -- (5.6,-1.1);
  \node[right] at (5.65,-0.55) {$v(L)$};
  \node[above] at (1.4,0.15) {$E = \SI{10}{GPa}$};
  \draw[dimto] (0,-1.85) -- (0.9,-1.85) node[right] {$x$};
  \draw[dim] (0,-2.45) -- (5,-2.45) node[midway, fill=white] {$L = \SI{0.5}{m}$};
  % cross-section
  \begin{scope}[shift={(-1.3,0)}]
    \draw[thick, fill=gray!12] (0,0) circle (0.35);
    \draw[dim] (-0.55,-0.35) -- (-0.55,0.35) node[midway, left] {\small$\SI{20}{mm}$};
  \end{scope}
\end{tikzpicture}
\end{minipage}%
\hfill
\begin{minipage}[c]{0.33\textwidth}
\textbf{Boundary conditions:}
\begin{enumerate}[label=\arabic*)]
  \item $v(0) = 0$
  \item $\dfrac{dv}{dx}(0) = 0$
  \item $M(L) = 0$
  \item $V(L) = P$
\end{enumerate}
\end{minipage}

\subsection*{Integrate four times}

\begin{align*}
  EI\,\frac{d^4 v}{dx^4} &= p(x) = 0 \\[4pt]
  \text{(i)} \qquad EI\,\frac{d^3 v}{dx^3} &= C_1 \\[4pt]
  \text{(ii)} \qquad EI\,\frac{d^2 v}{dx^2} &= C_1 x + C_2 \\[4pt]
  \text{(iii)} \qquad EI\,\frac{dv}{dx} &= \tfrac12 C_1 x^2 + C_2 x + C_3 \\[4pt]
  \text{(iv)} \qquad EI\,v(x) &= \tfrac16 C_1 x^3 + \tfrac12 C_2 x^2 + C_3 x + C_4
\end{align*}

\subsection*{Impose boundary conditions}

\begin{align*}
  \text{4)} \quad & V(L) = P = \left.\frac{d^3 v}{dx^3}\right|_L EI = C_1
     && \Rightarrow \quad \boxed{C_1 = P} \\[6pt]
  \text{3)} \quad & M(L) = 0 = \left.\frac{d^2 v}{dx^2}\right|_L EI = C_1 L + C_2 \\
  & \hphantom{M(L)}\;\;\, 0 = PL + C_2
     && \Rightarrow \quad \boxed{C_2 = -PL} \\[6pt]
  \text{2)} \quad & \left.\frac{dv}{dx}\right|_0 = 0 = \frac{1}{EI}\,C_3
     && \Rightarrow \quad \boxed{C_3 = 0} \\[6pt]
  \text{1)} \quad & v(0) = 0 = \frac{1}{EI}\,C_4
     && \Rightarrow \quad \boxed{C_4 = 0}
\end{align*}

\subsection*{Finally}

\[
  v(x) = \frac{1}{EI}\left[\tfrac16 P x^3 - \tfrac12 P x^2 L\right]
  \qquad\Rightarrow\qquad
  \fbox{\fbox{$\displaystyle v(x) = \frac{P x^2}{6EI}\,(x - 3L)$}}
\]

\noindent
$\rightarrow$ At the tip:
\[
  \fbox{\fbox{$\displaystyle v(L) = -\frac{PL^3}{3EI}$}}
\]

\noindent
For the squirrel problem, with $I = \frac{\pi}{64}(\SI{0.02}{m})^4 = \SI{7.85e-9}{m^4}$:
\[
  v_{\text{tip}} = -\frac{(\SI{5}{N})(\SI{0.5}{m})^3}{3\,(\SI{10e9}{Pa})(\SI{7.85e-9}{m^4})}
  \qquad\Rightarrow\qquad
  \boxed{v_{\text{tip}} = -\SI{2.65}{mm}}
\]

\end{document}
